\begin{abstract}
Let $\HH$ be a separable Hilbert space and let a prior be the terminal
law of an Ornstein--Uhlenbeck process perturbed by a learned,
Cameron--Martin-valued drift.  This construction is useful for networks
trained with score-matching objectives, but the drift in the defining
SDE is not assumed to be the reverse-time score of its own marginals.
Keeping that distinction explicit, we establish three measure-theoretic
facts for bounded linear observations with Gaussian noise.  First, any
such prior with a finite second moment gives a posterior that exists and
is locally Lipschitz in the data in Hellinger distance; a finite-horizon
Girsanov condition makes the terminal law equivalent to the reference
Gaussian.  Second, for a full-space lifted spectral
approximation, the posterior error satisfies
\[
  \dH(\muyN,\muy)
  \leq C_R\!\left(\sqrt{\varepsilon_N}+r_N\right),
  \qquad
  r_N^2=\sum_{j>N}\lambda_j ,
\]
where $\varepsilon_N$ is the pathwise Cameron--Martin drift error and
$(\lambda_j)$ are the reference-covariance eigenvalues.  The lift is
essential: a measure supported on a finite-dimensional subspace is
singular with respect to a nondegenerate infinite-dimensional prior.
Third, we state a conditional posterior-contraction criterion that
separates local prior mass from testing and identifiability.  It does
not imply a universal contraction exponent.  In particular, finite
dimensional observations identify an infinite-dimensional truth only
under an additional restriction on the prior support or an inverse
stability condition.
\end{abstract}
